% GENERATED FILE. Edit canonical lessons, then run npm run generate:books.
% canonical-source-hash: fnv1a64:df24aaa71227807a
% canonical-lessons: ZH-C110-waiyu, ZH-W110-wai, ZH-C110-gangqin, ZH-C110-tan, ZH-C110-huahua, ZH-C110-tiaowu, ZH-R110-review-105-110

\chapter{Things I Can Do}
\label{ch:zh-things-i-can-do}
\hlchaptermodality{110}

\begin{chapteropening}
\textbf{By the end of this chapter:} \emph{I can say I speak a foreign language, play the piano, paint and dance.}

Complete the last lesson of chapter 110: I can say I speak a foreign language, play the piano, paint and dance.
\end{chapteropening}

\section[\texorpdfstring{\zh{外语} (wàiyǔ)}{外语 (wàiyǔ)}]{\zh{外语} (wàiyǔ) — foreign language}

\label{lesson:ZH-C110-waiyu}



\begin{quote}
Before the new one: say the Mandarin for to charge a battery, then the Mandarin for to download.
\end{quote}

\subsection*{You'll want to know: \zh{外语}}
\textbf{\zh{外语}} — \emph{wàiyǔ} — \textquotedblleft{}foreign language\textquotedblright{}.

\begin{quote}
\textbf{\zh{我上网学外语}} \emph{wǒ shàngwǎng xué wàiyǔ} — I learn foreign languages online.
\end{quote}

The new character, in pieces: \zh{外} — on the left the three-stroke top half of \zh{多}; on the right a vertical with a dot beside it. The next lesson puts it in your hand.

\begin{cousinweb}
\textbf{\zh{外语}} (\emph{wàiyǔ}) is built from \zh{外} 'outside' and \zh{语} 'speech', both inherited from Middle Chinese, \emph{ngwajH} and \emph{ngjoX} (Guangyun, in Baxter's transcription): the speech from outside.
\end{cousinweb}

\begin{grammarlens}[title={What you've built}]
One more word for what I have and can do.
\end{grammarlens}

\subsection*{Your turn}
\begin{itemize}
\raggedright
  \item \emph{Say it:} \emph{wàiyǔ}
  \item \emph{Say it:} \emph{wàiyǔ}, once more
  \item \emph{Say it:} \emph{wǒ shàngwǎng xué wàiyǔ} — I learn foreign languages online.
  \item \emph{Recall:} say the Mandarin for to charge a battery, then the Mandarin for to download, then say \emph{wàiyǔ} again
\end{itemize}

\begin{tcolorbox}[breakable,colback=teal!4,colframe=teal!35!black,title={Before you move on}]
What does \zh{外语} mean? (\textquotedblleft{}Foreign language\textquotedblright{}.) Say it once more.
\end{tcolorbox}
\section[\texorpdfstring{\zh{外} (wài)}{外 (wài)}]{\zh{外} — the top of \zh{多}, then a vertical and a dot}

\label{lesson:ZH-W110-wai}



\begin{quote}
The word first, then the character inside it.

\begin{itemize}
\raggedright
  \item \emph{Recall:} say \emph{wàiyǔ} — \textbf{R1}, one lesson back, before it decays
\end{itemize}

\begin{quote}
\zh{多} \zh{语}
\end{quote}

\emph{Read it:} each one aloud, with nothing said first — \emph{duō}, \emph{yǔ}
\end{quote}

\subsection*{Script — the top of \zh{多}, then a vertical and a dot}
\hlblockfigure{\detokenize{figures/ZH-W110-wai-filmstrip.pdf}}{How \zh{外} is written, stroke by stroke}

\begin{quote}
\zh{外}
\end{quote}

\textbf{Two pieces, left before right.}

\textbf{The left piece, three strokes}, the top half of \zh{多}: a short left-falling stroke, a stroke across that turns and sweeps down to the left, and a dot inside.

\textbf{The right piece, two strokes}: a long vertical, and a dot to its right.

3 + 2 = five.

\subsection*{Your turn}
\begin{itemize}
\raggedright
  \item \emph{Write it:} \zh{外}, all five, with the model in front of you
  \item \emph{Cover:} write it again from memory, and count the pieces, not the strokes
\end{itemize}

\begin{tcolorbox}[breakable,colback=teal!4,colframe=teal!35!black,title={Before you move on}]
How many strokes in \zh{外}? (\textbf{Five}.) Where have you written its left piece? (\textbf{At the top of \zh{多}}.) What is on the right? (\textbf{A vertical and a dot}.)

Source: \href{https://raw.githubusercontent.com/chanind/hanzi-writer-data/68d10a4b21150cae5e1ebbd223eed289cf32d90c/data/\%E5\%A4\%96.json}{Hanzi Writer Data \zh{外}, pinned stroke paths and medians}.
\end{tcolorbox}
\section[gāngqín]{gāngqín — piano}

\label{lesson:ZH-C110-gangqin}



\begin{quote}
Before the new one: say the Mandarin for to download, then the Mandarin for foreign language.
\end{quote}

\subsection*{You'll want to know: gāngqín}
\emph{gāngqín} — \textquotedblleft{}piano\textquotedblright{}.

\begin{quote}
\emph{tā jiā yǒu gāngqín} — Her family has a piano.
\end{quote}

Its characters come later; for now it is a word for the ear.

\begin{cousinweb}
\emph{gāngqín} is built from gāng 'steel' and qín, a stringed instrument, both inherited from Middle Chinese, \emph{kang} and \emph{gim} (Guangyun, in Baxter's transcription); the dictionary describes the second as an instrument of five strings, later seven. The piano is a modern sense of the pair.
\end{cousinweb}

\begin{grammarlens}[title={What you've built}]
One more word for what I have and can do.
\end{grammarlens}

\subsection*{Your turn}
\begin{itemize}
\raggedright
  \item \emph{Say it:} \emph{gāngqín}
  \item \emph{Say it:} \emph{gāngqín}, once more
  \item \emph{Say it:} \emph{tā jiā yǒu gāngqín} — Her family has a piano.
  \item \emph{Recall:} say the Mandarin for to download, then the Mandarin for foreign language, then say \emph{gāngqín} again
\end{itemize}

\begin{tcolorbox}[breakable,colback=teal!4,colframe=teal!35!black,title={Before you move on}]
What does gāngqín mean? (\textquotedblleft{}Piano\textquotedblright{}.) Say it once more.
\end{tcolorbox}
\section[tán]{tán — to play an instrument; to pluck}

\label{lesson:ZH-C110-tan}



\begin{quote}
Before the new one: say the Mandarin for foreign language, then the Mandarin for piano.
\end{quote}

\subsection*{You'll want to know: tán}
\emph{tán} — \textquotedblleft{}to play an instrument; to pluck\textquotedblright{}.

\begin{quote}
\emph{wǒ huì tán gāngqín} — I can play the piano.
\end{quote}

Its characters come later; for now it is a word for the ear.

\begin{cousinweb}
\emph{tán} is inherited from Middle Chinese \emph{dan} (Guangyun, in Baxter's transcription), glossed 'to shoot' and named in a game of flicking pieces. Playing an instrument is not given there, so how the word came to that sense is not settled.
\end{cousinweb}

\begin{grammarlens}[title={What you've built}]
One more word for what I have and can do.
\end{grammarlens}

\subsection*{Your turn}
\begin{itemize}
\raggedright
  \item \emph{Say it:} \emph{tán}
  \item \emph{Say it:} \emph{tán}, once more
  \item \emph{Say it:} \emph{wǒ huì tán gāngqín} — I can play the piano.
  \item \emph{Recall:} say the Mandarin for foreign language, then the Mandarin for piano, then say \emph{tán} again
\end{itemize}

\begin{tcolorbox}[breakable,colback=teal!4,colframe=teal!35!black,title={Before you move on}]
What does tán mean? (\textquotedblleft{}To play an instrument; to pluck\textquotedblright{}.) Say it once more.
\end{tcolorbox}
\section[huàhuà]{huàhuà — to paint pictures; to draw}

\label{lesson:ZH-C110-huahua}



\begin{quote}
Before the new one: say the Mandarin for piano, then the Mandarin for to play an instrument.
\end{quote}

\subsection*{You'll want to know: huàhuà}
\emph{huàhuà} — \textquotedblleft{}to paint pictures; to draw\textquotedblright{}.

\begin{quote}
\emph{háizi xǐhuan huàhuà} — Children love painting.
\end{quote}

Its characters come later; for now it is a word for the ear.

\begin{cousinweb}
\emph{huàhuà} is built from huà 'to paint' and its object huà 'a picture', one word used twice, inherited from Middle Chinese \emph{hweaH} (Guangyun, in Baxter's transcription); the dictionary explains painting as laying the five colours on the likeness of things.
\end{cousinweb}

\begin{grammarlens}[title={What you've built}]
One more word for what I have and can do.
\end{grammarlens}

\subsection*{Your turn}
\begin{itemize}
\raggedright
  \item \emph{Say it:} \emph{huàhuà}
  \item \emph{Say it:} \emph{huàhuà}, once more
  \item \emph{Say it:} \emph{háizi xǐhuan huàhuà} — Children love painting.
  \item \emph{Recall:} say the Mandarin for piano, then the Mandarin for to play an instrument, then say \emph{huàhuà} again
\end{itemize}

\begin{tcolorbox}[breakable,colback=teal!4,colframe=teal!35!black,title={Before you move on}]
What does huàhuà mean? (\textquotedblleft{}To paint pictures; to draw\textquotedblright{}.) Say it once more.
\end{tcolorbox}
\section[tiàowǔ]{tiàowǔ — to dance}

\label{lesson:ZH-C110-tiaowu}



\begin{quote}
Before the new one: say the Mandarin for to play an instrument, then the Mandarin for to paint pictures.
\end{quote}

\subsection*{You'll want to know: tiàowǔ}
\emph{tiàowǔ} — \textquotedblleft{}to dance\textquotedblright{}.

\begin{quote}
\emph{nǐ xǐhuan tiàowǔ ma} — Do you like dancing?
\end{quote}

Its characters come later; for now it is a word for the ear.

\begin{cousinweb}
\emph{tiàowǔ} is built from tiào 'to leap' and wǔ 'dance', both inherited from Middle Chinese, \emph{dew} and \emph{mjuX} (Guangyun, in Baxter's transcription); the dictionary glosses the second as song and dance.
\end{cousinweb}

\begin{grammarlens}[title={What you've built}]
One more word for what I have and can do.
\end{grammarlens}

\subsection*{Your turn}
\begin{itemize}
\raggedright
  \item \emph{Say it:} \emph{tiàowǔ}
  \item \emph{Say it:} \emph{tiàowǔ}, once more
  \item \emph{Say it:} \emph{nǐ xǐhuan tiàowǔ ma} — Do you like dancing?
  \item \emph{Recall:} say the Mandarin for to play an instrument, then the Mandarin for to paint pictures, then say \emph{tiàowǔ} again
\end{itemize}

\begin{tcolorbox}[breakable,colback=teal!4,colframe=teal!35!black,title={Before you move on}]
What does tiàowǔ mean? (\textquotedblleft{}To dance\textquotedblright{}.) Say it once more.
\end{tcolorbox}
\section[(dialogue)]{Looking back: from In the Bathroom to Things I Can Do}

\label{lesson:ZH-R110-review-105-110}



\begin{quote}
Say the Mandarin for towel and toothbrush.
\end{quote}

\subsection*{Your turn}
\begin{itemize}
\raggedright
  \item \emph{Say it:} towel, toothbrush, toothpaste, soap, mirror
  \item \emph{Say it:} station, metro, bus, bicycle, to ride
  \item \emph{Say it:} to jump, to climb, to move house, to stop, to leave a place
  \item \emph{Say it:} to try on, to exchange, a pair of, size, a good fit
  \item \emph{Say it:} to go online, television, photo, to charge a battery, to download
  \item \emph{Say it:} foreign language, piano, to play an instrument, to paint pictures, to dance
  \item \emph{Read it:} \zh{站} \zh{试} \zh{网} \zh{外} — each one aloud, then a word that holds it
\end{itemize}

\begin{tcolorbox}[breakable,colback=teal!4,colframe=teal!35!black,title={Before you move on}]
Towel? (\emph{máojīn}.) Toothbrush? (\emph{yáshuā}.) Toothpaste? (\emph{yágāo}.) Soap? (\emph{féizào}.) Mirror? (\emph{jìngzi}.) Station? (\textbf{\zh{车站}}, \emph{chēzhàn}.) Metro? (\emph{dìtiě}.) Bus? (\emph{gōngjiāochē}.) Bicycle? (\emph{zìxíngchē}.) To ride? (\emph{qí}.) To jump? (\emph{tiào}.) To climb? (\emph{pá}.) To move house? (\emph{bān}.) To stop? (\emph{tíng}.) To leave a place? (\emph{líkāi}.) To try on? (\textbf{\zh{试}}, \emph{shì}.) To exchange? (\emph{huàn}.) A pair of? (\emph{shuāng}.) Size? (\textbf{\zh{大小}}, \emph{dàxiǎo}.) A good fit? (\emph{héshì}.) To go online? (\textbf{\zh{上网}}, \emph{shàngwǎng}.) Television? (\emph{diànshì}.) Photo? (\emph{zhàopiàn}.) To charge a battery? (\emph{chōngdiàn}.) To download? (\emph{xiàzài}.) Foreign language? (\textbf{\zh{外语}}, \emph{wàiyǔ}.) Piano? (\emph{gāngqín}.) To play an instrument? (\emph{tán}.) To paint pictures? (\emph{huàhuà}.) To dance? (\emph{tiàowǔ}.)
\end{tcolorbox}
